\part{Protected mode: unlocking the processor}

\chapter{Why real mode will not do}

\begin{goals}
  \item The three fatal limitations of real mode
  \item What the system gains by switching modes
  \item What is lost in the switch, and why timing matters
\end{goals}

\section{Three fatal limitations}

Real mode exists for compatibility with software from 1981. It has three problems that make it
useless for a genuine operating system:

\begin{description}
  \item[Only 1 MB of memory] With the \cod{segment$\times$16 + offset} scheme, the highest possible
    address is $\text{FFFF}\times16+\text{FFFF} = \text{10FFEF}$, barely over 1 MB. It does not
    matter that the machine has 16 GB: in real mode they do not exist.
  \item[16-bit registers] Everything is 16 bits. Pointers, counters, arithmetic. Manipulating a
    32-bit address requires constant juggling.
  \item[\textbf{No protection whatsoever}] And this is the decisive one. In real mode \emph{any}
    code can write to \emph{any} address, execute \emph{any} instruction, and switch off interrupts
    whenever it feels like it. There are no rings, no pages, no permissions. An operating system
    whose job is to isolate programs cannot run on a mode that offers no isolation mechanism.
\end{description}

\begin{concept}{What you gain}
32-bit \term{protected mode} provides exactly what is missing: 4 GB addressable, 32-bit registers,
the four privilege rings, segments with permissions, and --- most importantly --- the ability to
enable \hi{paging}, the foundation of virtual memory.
\end{concept}

\section{What you lose}

The switch is not free: on entering protected mode, \hi{BIOS services stop working}.
\cod{int 10h}, \cod{int 13h}, \cod{int 15h}: all of them are 16-bit code that assumes real mode.

Hence the ordering in the bootloader we saw in the previous part: print, read the disk and request
the memory map first (all through the BIOS), and only then switch modes. From that moment on,
anything you want from the hardware you will have to program yourself.

\begin{center}
\begin{tikzpicture}[font=\sffamily\scriptsize]
  \fill[amberlight,draw=amber,rounded corners=3pt] (0,0) rectangle (5.2,2.4);
  \node[font=\sffamily\bfseries,text=amber] at (2.6,2.1) {REAL MODE};
  \node[align=left,anchor=north west,font=\scriptsize] at (0.3,1.75)
    {\textcolor{green}{$\checkmark$} BIOS services\\
     \textcolor{red}{$\times$} Only 1 MB\\
     \textcolor{red}{$\times$} 16 bits\\
     \textcolor{red}{$\times$} No protection};

  \fill[greenlight,draw=green,rounded corners=3pt] (7.0,0) rectangle (12.2,2.4);
  \node[font=\sffamily\bfseries,text=green] at (9.6,2.1) {PROTECTED MODE};
  \node[align=left,anchor=north west,font=\scriptsize] at (7.3,1.75)
    {\textcolor{red}{$\times$} No BIOS services\\
     \textcolor{green}{$\checkmark$} 4 GB\\
     \textcolor{green}{$\checkmark$} 32 bits\\
     \textcolor{green}{$\checkmark$} Rings and paging};

  \draw[-{Stealth[length=3mm]},accent,line width=1.5pt] (5.4,1.2) -- (6.8,1.2);
  \node[font=\sffamily\tiny,text=accent,align=center] at (6.1,1.55) {one-way\\trip};
\end{tikzpicture}
\end{center}

% ============================================================
\chapter{The GDT: memory's driving licence}

\begin{goals}
  \item What a segment descriptor is and what information it carries
  \item How to decode a descriptor byte by byte
  \item What a selector is and how it relates to privilege rings
  \item Why we use a ``flat'' model and what that means
\end{goals}

\section{The problem it solves}

In protected mode, segment registers (\reg{CS}, \reg{DS}, \reg{SS}\ldots) no longer hold a number
to be multiplied by 16. They now hold a \term{selector}: an index into a table of
\term{descriptors}. And each descriptor states, for that region of memory:

\begin{itemize}
  \item where it starts (\emph{base}, 32 bits),
  \item how far it extends (\emph{limit}, 20 bits),
  \item what may be done with it (execute, read, write),
  \item \hi{from which privilege ring it may be used}.
\end{itemize}

That table is the \term{GDT} (Global Descriptor Table). It is the first thing that must be built,
because without it the processor refuses to enter protected mode.

\begin{analogy}
In real mode, a memory address is like saying ``house number 30 on the street that starts at
kilometre 7''. Anyone can name any house.

In protected mode you first show a badge (the selector). The badge refers to a record (the
descriptor) stating which area you may visit, whether you may enter or only look through the
window, and what clearance you need. The processor checks the record on \emph{every} access, in
hardware, at no cost. If you do not qualify, you do not get in.
\end{analogy}

\section{A descriptor, dissected}

Each descriptor occupies 8 bytes, with its fields chopped up and scattered in a way explicable only
by backwards compatibility with the 80286:

\begin{center}
\begin{tikzpicture}[font=\ttfamily\tiny,x=1.42cm,y=1cm]
  \foreach \i/\c/\t in {0/accentlight/{limit\\15--0},1/accentlight/{limit\\15--0},
                        2/purplelight/{base\\15--0},3/purplelight/{base\\15--0},
                        4/purplelight/{base\\23--16},5/amberlight/{access},
                        6/greenlight/{flags +\\limit 19--16},7/purplelight/{base\\31--24}} {
    \draw[draw=grayborder,fill=\c] (\i,0) rectangle (\i+0.94,1);
    \node[align=center,font=\sffamily\tiny] at (\i+0.47,0.5) {\t};
    \node[font=\ttfamily\tiny,text=gray] at (\i+0.47,-0.28) {byte \i};
  }
\end{tikzpicture}
\end{center}

The important field is the \term{access} byte. Its eight bits mean:

\begin{center}
\begin{tabularx}{\textwidth}{@{}c l X@{}}
\toprule
\textbf{Bit} & \textbf{Name} & \textbf{Meaning} \\
\midrule
7 & P    & Present. If 0, the descriptor is invalid and using it raises an exception \\
6--5 & DPL & \textbf{Privilege level}: 00 = ring 0 (kernel), 11 = ring 3 (user) \\
4 & S    & 1 = code or data segment; 0 = system segment (TSS, gates) \\
3 & E    & 1 = executable (code); 0 = data \\
2 & DC   & Direction / conforming \\
1 & RW   & For code: is it readable? For data: is it writable? \\
0 & A    & Accessed; set by the processor itself \\
\bottomrule
\end{tabularx}
\end{center}

\section{The lytlnyblOS GDT}

Our table has six entries. Here they are, with the access byte decoded:

\filetag{lytlnybl/kernel/kernel\_intro.asm}
\begin{lstlisting}[style=asm]
gdt_start:
gdt_null:
    dq 0                ; entry 0: mandatorily null

gdt_code:               ; selector 0x08 - kernel code
    dw 0xFFFF           ; limit low
    dw 0x0000           ; base low
    db 0x00             ; base middle
    db 0x9A             ; access: present, ring 0, code, readable
    db 0xCF             ; flags: 4K granularity + 32-bit, limit high
    db 0x00             ; base high

gdt_data:               ; selector 0x10 - kernel data
    dw 0xFFFF
    dw 0x0000
    db 0x00
    db 0x92             ; access: present, ring 0, data, writable
    db 0xCF
    db 0x00

gdt_user_code:          ; selector 0x18 - user code
    dw 0xFFFF
    dw 0x0000
    db 0x00
    db 0xFA             ; access: present, RING 3, code, readable
    db 0xCF
    db 0x00

gdt_user_data:          ; selector 0x20 - user data
    dw 0xFFFF
    dw 0x0000
    db 0x00
    db 0xF2             ; access: present, RING 3, data, writable
    db 0xCF
    db 0x00

gdt_tss:                ; selector 0x28 - filled in later, from C
    dw 0
    dw 0
    db 0
    db 0
    db 0
    db 0
gdt_end:
\end{lstlisting}

\subsection{Decoding \hexa{9A} and \hexa{FA}}

These two bytes are the difference between the kernel and a user program. Bit by bit:

\begin{center}
\begin{tabularx}{\textwidth}{@{}l c c c c c c c c X@{}}
\toprule
 & \textbf{P} & \multicolumn{2}{c}{\textbf{DPL}} & \textbf{S} & \textbf{E} & \textbf{DC} & \textbf{RW} & \textbf{A} & \\
\midrule
\hexa{9A} = \cod{10011010} & 1 & 0 & 0 & 1 & 1 & 0 & 1 & 0 & code, \textbf{ring 0} \\
\hexa{FA} = \cod{11111010} & 1 & 1 & 1 & 1 & 1 & 0 & 1 & 0 & code, \textbf{ring 3} \\
\hexa{92} = \cod{10010010} & 1 & 0 & 0 & 1 & 0 & 0 & 1 & 0 & data, ring 0 \\
\hexa{F2} = \cod{11110010} & 1 & 1 & 1 & 1 & 0 & 0 & 1 & 0 & data, ring 3 \\
\bottomrule
\end{tabularx}
\end{center}

\hi{Those two DPL bits are the boundary between the two worlds of Chapter 7.} All of the
system's isolation rests, literally, on them.

\subsection{Base 0, maximum limit: the flat model}

Notice that the four useful descriptors all have base \hexa{00000000} and limit \hexa{FFFFF} with
4 KB granularity, which covers the full 4 GB.

That means \emph{every segment is the same segment}: all of memory. It is called the \term{flat
model}, and its advantage is enormous: the logical address you write in C is the linear address,
with no strange translation. Segments stop being a way to separate memory and are left doing
\hi{one job only}: carrying the privilege bits.

\begin{concept}{Division of labour in modern x86}
Since the 386, memory isolation is done with \hi{paging}, not segmentation. Segments are left
flat and used only for privilege level. That is exactly what Linux does, and what lytlnyblOS does.
Segmentation survives as a mandatory piece of paperwork you must fill in before you can use the
thing that actually matters.
\end{concept}

\section{Selectors}

A selector is a 16-bit number with three fields:

\begin{center}
\begin{tikzpicture}[font=\sffamily\scriptsize,x=1cm]
  \draw[draw=grayborder,fill=accentlight] (0,0) rectangle (6.5,0.85);
  \node at (3.25,0.42) {index into the table (13 bits)};
  \draw[draw=grayborder,fill=graylight] (6.5,0) rectangle (7.7,0.85);
  \node at (7.1,0.42) {TI};
  \draw[draw=grayborder,fill=redlight] (7.7,0) rectangle (9.5,0.85);
  \node at (8.6,0.42) {RPL};
  \node[font=\ttfamily\tiny,text=gray] at (0.4,-0.3) {bit 15};
  \node[font=\ttfamily\tiny,text=gray] at (9.1,-0.3) {bit 0};
\end{tikzpicture}
\end{center}

Because the low two bits are the RPL, the index is shifted three positions up. Hence selectors go
in steps of eight:

\begin{center}
\begin{tabularx}{0.95\textwidth}{@{}l l l X@{}}
\toprule
\textbf{Selector} & \textbf{Index} & \textbf{RPL} & \textbf{Points at} \\
\midrule
\hexa{00} & 0 & 0 & null descriptor \\
\hexa{08} & 1 & 0 & kernel code \\
\hexa{10} & 2 & 0 & kernel data \\
\hexa{18} \cod{| 3} = \hexa{1B} & 3 & \textbf{3} & user code \\
\hexa{20} \cod{| 3} = \hexa{23} & 4 & \textbf{3} & user data \\
\hexa{28} & 5 & 0 & the TSS \\
\bottomrule
\end{tabularx}
\end{center}

Which is why creating a user process looks like this:

\filetag{lytlnybl/tasks/procman.c}
\begin{lstlisting}[style=c]
void create_uprocess(process_t* new_process) {
    new_process->regs.cs = 0x18 | 3;   /* user code, RPL = 3 */
    new_process->regs.ds = 0x20 | 3;   /* user data, RPL = 3 */
    new_process->regs.ss = 0x20 | 3;
    /* ... */
}
\end{lstlisting}

\exercise{What would happen if a user process were created with \cod{cs = 0x18} without the
\cod{| 3}? Think about what the processor checks when loading \reg{CS}. (Hint: one rule says the
RPL may not be more privileged than the descriptor's DPL, and another forbids ``dropping''
privilege through a plain \cod{iret}.)}

% ============================================================
\begin{recipe}{reading the live GDT}{ch16-gdt}
Asks the processor for its own GDT with \cod{SGDT}, walks the six descriptors, decodes each access
byte into P / DPL / S / E, and prints which selectors the kernel is currently running with.

It surfaces something the source alone cannot tell you: \filepath{kernel\_intro.asm} writes
\hexa{92} for the kernel data segment, but the live table reads back \hexa{93}. Bit 0 is
\emph{Accessed}, and the processor sets it the first time the descriptor is loaded into a segment
register. The TSS likewise reads \hexa{8B} rather than \hexa{89}, because \cod{ltr} sets bit 1,
\emph{Busy}. The hardware writes into your GDT.
\end{recipe}

% ============================================================
\chapter{The jump to 32 bits}

\begin{goals}
  \item The exact sequence for entering protected mode
  \item Why a far jump is needed immediately afterwards
  \item What the kernel stack is and where it lives
\end{goals}

\section{The sequence}

\filetag{lytlnybl/kernel/kernel\_intro.asm}
\begin{lstlisting}[style=asm]
enter_protected:
    cli                     ; 1. disable interrupts
    lgdt [gdtr - start]     ; 2. load the GDT
    mov eax, cr0
    or eax, 1               ; 3. set the PE bit in CR0
    mov cr0, eax

    CODE_SEG equ gdt_code - gdt_start
    jmp CODE_SEG:p_mode_main ; 4. far jump -> reload CS
[bits 32]
p_mode_main:
    mov ax, 10h             ; 5. reload the data segments
    mov ds, ax
    mov es, ax
    mov fs, ax
    mov gs, ax
    mov ss, ax
    mov esp, kernel_stack_top ; 6. set up the stack

    mov byte [0xB8000], 'P'   ; 7. sign of life
    mov byte [0xB8001], 0x02

    extern kernel_main
    call kernel_main          ; 8. enter C
\end{lstlisting}

\subsection{Step 1: \cod{cli}}

Disables interrupts. It is essential: for the next few microseconds the processor is in an
inconsistent state, with a new GDT but old segments. If an interrupt arrived, it would jump into
the \emph{real mode} vector table, which no longer makes sense, and the machine would hang.

Interrupts are not re-enabled until the end of \cod{idt\_init()}, already in C, once a valid 32-bit
interrupt table exists.

\subsection{Step 2: \cod{lgdt}}

Loads the \reg{GDTR} register with the table's address and size. The operand format is peculiar:

\begin{lstlisting}[style=asm]
gdtr:
    dw gdt_end - gdt_start - 1   ; limit: size MINUS ONE
    dd gdt_start                 ; base address, 32 bits
\end{lstlisting}

\begin{warn}{The minus one}
The limit is ``the last valid byte'', not ``the size''. With 6 descriptors of 8 bytes that is 48
bytes, so the limit is 47. This is a classic bug: put 48 and the last descriptor falls partly out of
range, and using it raises a general protection fault. The same convention applies to the IDT.
\end{warn}

\subsection{Step 3: the PE bit}

\cod{or eax, 1} sets bit 0 of \reg{CR0}. At that instant the processor is \emph{already} in
protected mode. But there is a subtlety.

\subsection{Step 4: why the far jump is required}

The processor internally caches (in a ``descriptor cache'') the information for the current
segment, and that cache still holds the real-mode interpretation of \reg{CS}. Until \reg{CS} is
reloaded, the processor keeps executing code as if it were 16-bit.

And \reg{CS} \emph{cannot be modified with a \cod{mov}}. The only ways to change it are a far jump,
a far call, or an \cod{iret}. Hence:

\begin{lstlisting}[style=asm]
jmp CODE_SEG:p_mode_main
\end{lstlisting}

That jump is the exact moment the processor begins executing 32-bit instructions. The \cod{[bits
32]} directive right after it generates no code: it tells NASM to assemble for 32 bits from there.

\begin{concept}{The instant of the switch}
Three instructions mark the transition: \cod{mov cr0, eax} enables the mode, the far \cod{jmp}
reloads \reg{CS}, and the \cod{mov}s reload the data segments. Until all three are done, the
processor sits in an intermediate state where almost nothing behaves as you expect. That is why the
whole block runs with interrupts off and calls nothing.
\end{concept}

\subsection{Step 6: the stack}

\begin{lstlisting}[style=asm]
section .bss
align 16
kernel_stack_bottom:
    resb 0x4000          ; 16 KB of stack
kernel_stack_top:
\end{lstlisting}

\cod{resb} reserves space without initialising it, in the \cod{.bss} section. Notice that
\cod{kernel\_stack\_top} sits \emph{after} the 16 KB: since the stack grows downwards, \reg{ESP}
must point at the end of the block.

\begin{danger}{16 KB is not much}
That is the entire kernel stack. Every function call consumes tens or hundreds of bytes; every
local array, whatever it occupies. If the kernel overruns it, \reg{ESP} drops below
\cod{kernel\_stack\_bottom} and starts clobbering whatever is there --- usually the kernel's own
global variables. There is no check whatsoever: nobody warns you.

Part XIII shows a real case where the system rebooted in a loop for exactly this reason, and how it
was diagnosed.
\end{danger}

% ============================================================
\chapter{Entering C}

\begin{goals}
  \item What compiling ``freestanding'' means and what disappears
  \item What a linker script does and why it is mandatory here
  \item How the \cod{.text}, \cod{.data} and \cod{.bss} sections relate
\end{goals}

\section{C without a standard library}

Writing the kernel in C means giving up nearly everything you normally take for granted. The
compiler's \cod{-ffreestanding} option says so:

\filetag{lytlnybl/Makefile}
\begin{lstlisting}[style=mk]
CC = i686-elf-gcc
CFLAGS = -g -m32 -Iinclude -ffreestanding
\end{lstlisting}

\begin{center}
\begin{tabularx}{\textwidth}{@{}l X@{}}
\toprule
\textbf{Flag} & \textbf{What it does} \\
\midrule
\cod{-m32} & Generate 32-bit code \\
\cod{-ffreestanding} & Do not assume a standard library or a \cod{main()} exists; do not emit implicit calls to library functions \\
\cod{-Iinclude} & Look for headers in \filepath{include/} \\
\cod{-g} & Include debug information for GDB \\
\bottomrule
\end{tabularx}
\end{center}

What \emph{disappears}: \cod{printf}, \cod{malloc}, \cod{strlen}, \cod{FILE}, all of it. If the
kernel wants to copy memory, somebody has to write \cod{memcpy}. And somebody does:

\filetag{lytlnybl/kernel/kernel\_utils.c}
\begin{lstlisting}[style=c]
void* memcpy(void* dest, const void* src, size_t len) {
  uint8_t* d = (uint8_t*)dest;
  const uint8_t* s = (const uint8_t*)src;
  for (size_t i = 0; i < len; i++)
    d[i] = s[i];
  return dest;
}
\end{lstlisting}

\begin{warn}{What you can still use}
The headers \cod{<stdint.h>}, \cod{<stddef.h>} and \cod{<stdbool.h>} \emph{are} available even in
freestanding mode, because they define no functions: only types (\cod{uint32\_t}, \cod{size\_t},
\cod{bool}). You will see them included throughout the project.
\end{warn}

\section{The linker script}

The compiler produces chunks of code. The linker joins them and decides what address each one goes
to. Normally that decision is made by default, but here \hi{the address matters}: the
bootloader jumps to \hexa{9000} and the start of the kernel must be there.

\filetag{lytlnybl/kernel/linker.ld}
\begin{lstlisting}[style=txt]
ENTRY(start)

SECTIONS
{
    . = 0x9000;          /* the location counter starts here */

    _kernel_start = .;   /* remember where we start */

    .text : { *(.text) } /* executable code */
    .data : { *(.data) } /* initialised globals */
    .bss  : { *(.bss)  } /* zeroed globals */

    _kernel_end = .;     /* remember where we end */
}
\end{lstlisting}

The dot \cod{.} is the \term{location counter}: the address where the next thing will be placed.
Assigning \hexa{9000} to it and then listing sections is, literally, placing pieces in memory.

\subsection{The three sections}

\begin{center}
\begin{tabularx}{\textwidth}{@{}l l X@{}}
\toprule
\textbf{Section} & \textbf{In the file?} & \textbf{Contains} \\
\midrule
\cod{.text} & yes & Machine code. Read-only in serious systems \\
\cod{.data} & yes & Globals with an initial value: \cod{int x = 42;} \\
\cod{.bss}  & \textbf{no} & Globals with no value or zero: \cod{int y;}. Only space is reserved \\
\bottomrule
\end{tabularx}
\end{center}

\begin{concept}{Why \cod{.bss} takes no file space}
If you have a global 16 KB array of zeros, it would be absurd to store 16 KB of zeros on disk. The
linker notes ``16 KB of blank goes here'' and whoever loads the program is responsible for zeroing
it.

\hi{And here is a real trap in this project}: the thing that loads the kernel is the
bootloader, which zeroes nothing; it simply copies sectors. It works because RAM happens to be zero
at boot. It is a detail a serious system would have to handle explicitly.
\end{concept}

\subsection{The \cod{\_kernel\_start} and \cod{\_kernel\_end} symbols}

These two labels are not variables: they are \emph{addresses} computed by the linker. From C they
are declared like this:

\filetag{lytlnybl/memory/pmm.c}
\begin{lstlisting}[style=c]
extern char _kernel_start;
extern char _kernel_end;

/* and used by taking their ADDRESS, not their value */
uint32_t kernel_frame_index = (uint32_t)&_kernel_start / FRAME_SIZE;
uint32_t kernel_frame_end   = (uint32_t)&_kernel_end   / FRAME_SIZE;
\end{lstlisting}

They are declared as \cod{char} and used with \cod{\&} because what matters is \emph{where they
are}, not what they contain. Thanks to them the memory manager can mark exactly the pages the
kernel occupies as used, with no magic numbers.

\begin{danger}{A latent bug hiding in this division}
Look at those two lines again. \cod{\_kernel\_end} almost never lands exactly on a page boundary:
if it is \hexa{13468}, integer division gives 19, and the loop protects pages 9 through
\textbf{18}. \hi{Page 19 --- holding the last 1,128 bytes of the kernel --- is left
unprotected}, and the memory manager may hand it to anyone.

If important global variables live in that page, the system overwrites itself. Today the project
is lucky: the critical variables fall below the boundary. A few extra lines of code and that luck
runs out. The fix is one line --- round up --- and it is applied in the current source. Part XIII
gives the complete diagnosis, with the
processor trace.
\end{danger}

\section{And finally, C}

\begin{lstlisting}[style=asm]
extern kernel_main
call kernel_main
\end{lstlisting}

\cod{extern} tells NASM that the linker will resolve this symbol. \cod{call} pushes the return
address and jumps. From that instruction on, we are executing compiled C, with a valid stack, 32
bits and protected mode.

\filetag{lytlnybl/kernel/kernel\_main.c}
\begin{lstlisting}[style=c]
void kernel_main(void)
{
    volatile char* vga = (volatile char*)0xB8000;
    vga[0] = 'C';       /* sign of life: we reached C */
    vga[1] = 0x02;      /* attribute: green on black */

    vga_text_init(&terminal);
    vga_text_writeline(&terminal, "Welcome to the lytlnyblOS!");

    idt_init();
    init_pmm();
    init_vmm();
    init_heap();
    init_procman();
    timer_init(100);
    keyboard_init();
    init_tss();
    fs_manager_init();
    /* ... */
}
\end{lstlisting}

That initialisation order is not arbitrary: it is exactly the dependency chain from Chapter 6.
Interrupts before anything else, because if something fails we want to see the message. Physical
memory before virtual, because virtual needs to request pages from it. The heap after virtual,
because it needs to map. Processes after the heap, because they allocate structures with
\cod{kmalloc}.

\begin{concept}{Why \cod{volatile}}
\cod{volatile} tells the compiler: ``this pointer refers to something that can change for reasons
you cannot see; do not optimise these accesses away''. Without it, an optimising compiler could
delete the write to \hexa{B8000} entirely, reasoning that nobody reads that variable afterwards. In
code that talks to hardware, \cod{volatile} is not optional.
\end{concept}

\begin{tryit}
Boot with \cod{make run}. If a green \cod{P} appears in the top-left corner and then a green
\cod{C}, you have passed both milestones: protected mode and entry into C. They are the two signs of
life the code prints by hand, before any video driver exists.
\end{tryit}

\begin{summary}
The processor is in protected mode, 32-bit, with a GDT holding four flat segments plus the TSS, its
own stack, and executing C. We now have a decent programming environment. Next: being able to see
what is happening.
\end{summary}
